\documentclass[../../../include/open-logic-section]{subfiles}

\begin{document}

\olfileid{sth}{story}{extensionality}
\olsection{گسترش‌مندی}

نخستین نکته این است که مجموعه‌ها با !!{element}هایشان تشخص می‌یابند.
به بیان دقیق‌تر:

\begin{axiom}[گسترش‌مندی]
اگر مجموعه‌های~$A$ و~$B$، !!{element}های یکسانی داشته باشند، آنگاه
$A$ و~$B$ یک مجموعه‌اند.
\[
  \lforall[A][\lforall[B][(\lforall[x][(x \in A \liff x \in B)] \lif
  \eq[A][B])]]
\]
\end{axiom}

این اصل را در سراسر~\olref[sfr][][]{part} مفروض گرفتیم، اما ارزش تکرار
دارد. اصل موضوع گسترش‌مندی بیانگر این اندیشهٔ پایه است که یک مجموعه با
!!{element}هایش تعیین می‌شود. (بنابراین می‌توان مجموعه‌ها را با
\emph{مفاهیم} مقایسه کرد، که اشیای دقیقاً یکسانی ممکن است ذیل مفاهیم
گوناگون بسیاری قرار گیرند.)

چرا باید این اصل را بپذیریم؟ می‌توان به‌وجهی معقول گفت هرگونه انکار
گسترش‌مندی تصمیمی است برای کنارگذاشتن هر آنچه حتی بتوان
\emph{نظریهٔ مجموعه‌ها} نامید. نظریهٔ مجموعه‌ها چیزی بیش یا کم‌تر از
نظریهٔ گردایه‌های گسترشی نیست.

اما چالش واقعی در~\olref[sth][][]{part} وضع اصولی است که بگویند
\emph{کدام مجموعه‌ها وجود دارند}. و معلوم می‌شود تنها پاسخ حقیقتاً
«بدیهی» به این پرسش، به‌طور اثبات‌پذیر نادرست است.

\end{document}
